% ABSTRACT--------------------------------------------------------------------------------

\begin{resumo}[ABSTRACT]
\begin{SingleSpacing}

\imprimirautorcitacao. \imprimirtitleabstract. \imprimirdata. \pageref{LastPage} f. \imprimirprojeto\ -- \imprimirprograma, \imprimirinstituicao. \imprimirlocal, \imprimirdata.\\

Formal employment in Information Technology (IT) in Brazil has grown rapidly over the past two decades, yet the available analyses of the sector typically rely on closed aggregates or short time windows, and cannot be reproduced. This work built a complete analytical platform over the microdata of the Annual Social Information Report (RAIS) and the General Register of Employed and Unemployed Workers (CAGED), covering 2007 to 2026, organized as a Medallion architecture (Bronze, Silver and Gold) that runs entirely in the cloud: the data lake lives in public Hugging Face Hub repositories, and processing is performed by GitHub Actions workflows chained into a dependency graph, which periodically checks the public server of the Ministry of Labour and Employment and reprocesses only the years with new data. Development followed CRISP-DM, and its iterative nature was decisive: findings in the data led to new sources, a revised scope and reformulated questions. The Apache Parquet files are queried in place by the in-process engine DuckDB, with no prior import, and presented through a web dashboard built with React and published as a static site, in which DuckDB, compiled to WebAssembly, runs the queries in the visitor's own browser. A total of 661~GB of raw microdata in text form was ingested and stored as 56.4~GB of Parquet; the technology scope, defined as the union of economic activity (CNAE) and occupation (CBO), yielded 23.8 million RAIS employment records and 12.1 million CAGED movements. Reproducibility was verified by recounting the scope in the raw data for every year of RAIS and CAGED, which exposed and corrected a duplicate in 2007 and ended with zero difference, and the totals of the occupation lens fell within about 1\% of those in a public panel published by the Ministry. The results show that the stock of formal IT jobs rose from 560,040 in 2007 to 1,361,231 in 2025, a 143\% increase, while median pay measured in minimum wages changed little (3.24 to 3.41); that 58\% of IT professionals work outside IT companies, with median pay close to that of those inside them (4.13 against 4.31 minimum wages); that the gender wage gap fell from 19.0\% to 7.8\% but remained entirely unexplained by observable characteristics; and that a nowcasting estimator for the employment stock, validated by holding out each year from 2014 to 2025, was off by 2.1\% on average, against 5.1\% for the naive estimator. For the monthly balance, the SARIMA model selected through rolling-origin validation outperformed the seasonal naive baseline by 33\% in mean absolute error, against 10.5\% for the best Prophet configuration. Because the data lake itself is public, the processed layers and the aggregates are at once the source of the cloud-deployed dashboard and open datasets that third parties can reuse.\\

\textbf{Keywords}: IT labor market; data engineering; RAIS; CAGED; open data.

\end{SingleSpacing}
\end{resumo}
